\section{Introduction}

Entity authentication forms the cornerstone of modern cybersecurity, serving as the first line of defense in protecting digital assets and systems. As the digital landscape continues to evolve, understanding the strengths and vulnerabilities of authentication mechanisms becomes increasingly critical for security professionals.

\subsection{Background and Motivation}

Authentication systems rely on multiple factors: \textit{what you know} (passwords), \textit{what you have} (tokens), \textit{who you are} (biometrics), and \textit{where you are} (location). This study focuses primarily on password-based authentication and token-based systems, examining their implementation across different platforms and their resistance to various attack vectors.

Password authentication has been the dominant authentication method for decades, implemented at both the system level (Linux, Windows) and web application level (HTML forms, HTTP authentication). However, the security of password-based systems depends heavily on the underlying cryptographic implementation, particularly the password verification data (PVD) storage mechanisms.

\subsection{Security Challenges}

Modern authentication systems face several critical challenges:

\begin{enumerate}
    \item \textbf{Legacy Systems}: Older systems often employ weak hashing algorithms without proper salt or iteration counts
    \item \textbf{User Behavior}: Users frequently choose weak, predictable passwords that are vulnerable to dictionary attacks
    \item \textbf{Implementation Flaws}: Web applications may lack proper rate limiting or account lockout mechanisms
    \item \textbf{Attack Evolution}: Attackers have access to powerful tools and massive wordlists, making brute-force attacks increasingly effective
\end{enumerate}

\subsection{Research Objectives}

This assessment aims to achieve several key learning objectives:

\begin{enumerate}
    \item Understand the roles of salt and iteration count in mitigating offline dictionary attacks
    \item Differentiate between offline dictionary attacks and online password guessing
    \item Analyze token authentication with private keys
    \item Develop the ability to quickly identify the weakest authentication links in systems
\end{enumerate}

\subsection{Scope and Limitations}

This study examines three distinct computing environments representing different eras of password security: Windows 2003 Server (legacy LM hashes), Windows 7 Enterprise (NTLM hashes), and Ubuntu 14.04 (modern salted hashes). Additionally, we analyze web-based authentication mechanisms including HTML form authentication, HTTP Basic authentication, and HTTP Digest authentication.

All testing was conducted in controlled laboratory environments with proper authorization, following JMU's Honor Code and ethical hacking guidelines.

\subsection{Document Organization}

This report is structured to provide comprehensive coverage of all assessment activities:
\begin{itemize}
    \item \textbf{Methods}: Tools, techniques, and experimental setup
    \item \textbf{Results}: Task-by-task findings aligned with project deliverables
    \item \textbf{Discussion}: Analysis of security implications and comparative assessment
    \item \textbf{Conclusion}: Key findings and recommendations
    \item \textbf{Appendix}: Complete evidence documentation and supporting materials
\end{itemize}